\chapter{Hardware-/Software-Allokation und ROS2-Mapping}

\section{Allokationsprinzip}

Die Hardware-/Software-Allokation wird aus Funktionen abgeleitet. Erst wenn die Funktion, ihr Nachweis und ihre Schnittstellen verstanden sind, werden konkrete Komponenten festgelegt. Dies ist besonders wichtig fuer Sensorik, Rechnerplattform, Motorregelung und Safety.

\section{Vorlaeufige Allokation}

\begin{longtable}{L{0.18\textwidth} L{0.26\textwidth} L{0.28\textwidth} L{0.14\textwidth}}
\toprule
Funktion & Hardware-Anteil & Software-Anteil & Status \\
\midrule
FUN-MAP-001 & Onboard-Speicher, Rechner & Kartenserver, Kartenverwaltung, Versionierung & MVP \\
FUN-LOC-001 & Sensorik, Encoder, IMU, Rechner & Lokalisierungsalgorithmus, Transformationsbaum & MVP \\
FUN-PER-001 & LiDAR/Kamera/weitere Sensorik & Sensortreiber, Filterung, Hindernisextraktion & MVP \\
FUN-NAV-001 & Rechner & globale Planung, lokale Planung, Recovery-Verhalten & MVP \\
FUN-MOT-001 & Motoren, Motorcontroller, Encoder & hardwarenahe Regelung, \ros-Anbindung & MVP \\
FUN-SAFE-001 & Sensorik, Not-Halt, ggf. Safety-Hardware & Ueberwachung, Stopp-Logik, Fehlerstatus & MVP reduziert \\
FUN-TRN-001 & Mechanik der Aufnahme, Aktoren, Sensorik & Handhabungslogik, Verriegelungsstatus & Zielsystem \\
\bottomrule
\end{longtable}

\section{ROS2-Mapping}

Das Lehrbuch beschreibt das Mapping von \sysml-Elementen auf \ros-Elemente als Architekturentscheidung, nicht als starre Uebersetzungsregel \cite[Kap.~16]{becker2026sysml}. Ein \sysml-Element kann mehrere Nodes beeinflussen, und ein Node kann mehrere modellierte Verantwortlichkeiten buendeln.

\begin{longtable}{L{0.20\textwidth} L{0.20\textwidth} L{0.22\textwidth} L{0.22\textwidth}}
\toprule
\sysml-Element & \ros-Element & Schnittstellenidee & Verknuepfte Anforderungen \\
\midrule
MissionSystem & \texttt{mission\_manager\_node} & Auftrag, Ziel, Mission State & REQ-MIS-001, REQ-MIS-002, REQ-STA-001 \\
UserInterfaceSystem, CommunicationSystem & \texttt{web\_interface\_node} & Auftraege und Bedienbefehle von der Web-App, Status und Diagnose zur Web-App & REQ-UI-001, REQ-UI-002 \\
NavigationSystem & \texttt{navigation\_node}, z.\,B. auf Basis Nav2 & Zielaktion, Pfad, Bewegungsbefehl & REQ-NAV-001 bis REQ-NAV-005, REQ-PERF-001 \\
NavigationSystem (Kartenverwaltung) & \texttt{map\_management\_node} & Karte bereitstellen, Kartenupdate-Kandidaten & REQ-MAP-001, REQ-MAP-003 \\
NavigationSystem (Lokalisierung) & \texttt{localization\_node} & Pose, Transformationsdaten & REQ-LOC-001 \\
PerceptionSystem & \texttt{perception\_node}, ggf. aufgeteilt in Sensortreiber und Hindernisverarbeitung & Sensordaten, Hindernisse, lokale Karte & REQ-MAP-002, REQ-SAFE-002 \\
MotionSystem & \texttt{motion\_controller\_node}, ggf. mit MCU-Firmware & \texttt{cmd\_vel}, Odometrie, Motorstatus & REQ-NAV-001 \\
SafetySystem & \texttt{safety\_supervisor\_node}, spaeter plus Hardware-Stopp & Safety-Zustand, Stoppgrund & REQ-SAFE-001, REQ-SAFE-002 \\
EnergySystem & \texttt{energy\_monitor\_node} & Batterie- und Warnstatus & REQ-PWR-001, REQ-PWR-002 \\
DiagnosticsSystem & \texttt{diagnostics\_node} & Diagnosemeldungen & REQ-DIAG-001 \\
\bottomrule
\end{longtable}

\section{Vorlaeufige ROS2-Schnittstellen}

\begin{longtable}{L{0.18\textwidth} L{0.20\textwidth} L{0.16\textwidth} L{0.30\textwidth}}
\toprule
Name & Typ & Richtung & Zweck \\
\midrule
\texttt{/goal\_pose} & Topic/Action offen & Mission zu Navigation & Zielvorgabe fuer Navigation \\
\texttt{/cmd\_vel} & Topic & Navigation zu Motion & Geschwindigkeitsvorgabe \\
\texttt{/odom} & Topic & Motion zu Navigation & Odometriedaten \\
\texttt{/map} & Topic/Service & Kartenverwaltung zu Navigation/Lokalisierung & Grundkarte \\
\texttt{/local\_obstacles} & Topic & Perception zu Navigation & erkannte Hindernisse \\
\texttt{/mission\_state} & Topic & Mission zu UI/Logging & Missionsstatus \\
\texttt{/safety\_state} & Topic & Safety zu Mission/Motion & Safety-Zustand und Stoppgrund \\
\texttt{/battery\_state} & Topic & Energy zu Mission/UI & Energiezustand \\
\texttt{/transport\_order} & Topic & Web-Interface zu Mission & Transportauftrag \\
\texttt{/operator\_command} & Topic & Web-Interface zu Mission & Status, Stopp, Fortsetzung, Rueckkehr \\
\texttt{/navigation\_result} & Topic & Navigation zu Mission & Ziel erreicht, blockiert, nicht erreichbar \\
\texttt{/robot\_pose} & Topic & Lokalisierung zu Navigation & geschaetzte Pose \\
\texttt{/map\_update\_candidate} & Topic & Kartenverwaltung & vorgeschlagene dauerhafte Kartenaenderung \\
\texttt{/diagnostics} & Topic & Diagnose zu Web-Interface & Diagnosemeldungen \\
\bottomrule
\end{longtable}

\section{Noch nicht entschiedene Alternativen}

\begin{itemize}
  \item Rechnerplattform: Raspberry Pi, Jetson, x86 oder Kombination mit MCU.
  \item Sensorik: 2D-LiDAR, Tiefenkamera, Stereo, Ultraschall oder kombinierte Sensorik.
  \item Motorregelung: reine Motorcontroller-Loesung, MCU-Firmware oder integrierte Servoregler.
  \item Navigation: ROS2 Nav2 als Referenz, Eigenanteile fuer lokale Hindernislogik, alternative SLAM-/Lokalisierungspakete.
  \item Simulation: Gazebo, Isaac Sim oder reduzierter Testsimulator fuer fruehe MVP-Tests.
\end{itemize}
